The v7 benchmark yields three main insights for multimodal GBM prognosis.

\textbf{Cross-scale priors are biologically meaningful even when predictive gains are small.} The single-cell branch identifies recurrent GBM marker genes and compresses them into six transferable programs. This is an interpretable bridge between cell-state resolution and patient-level prediction, and it is the clearest conceptual flashpoint of the study. Rather than treating scRNA-seq as a separate descriptive add-on, the prior-guided model turns recurrent markers into a concrete attention constraint. This makes the architecture scientifically legible even when its current predictive gain is modest.

\textbf{The strongest empirical gain came from methylation, not from the current scRNA transfer.} The CPTAC ablation makes this point unambiguous. Adding methylation to RNA+clinical features substantially improved 3-year discrimination, especially for the SVM branch, whereas adding the current six scRNA program scores changed AUC only marginally. This suggests that bulk methylation carries a stronger immediately usable survival signal in this cohort, while scRNA-derived priors remain limited by the fact that they were learned from only 17 source cases.

\textbf{Strict external validation is more informative than pooled cross-validation.} The compact transformer achieved the best 1-year external AUC and LASSO logistic achieved the best 3-year external AUC, but all models remained in a relatively narrow range. The non-significant external Kaplan--Meier split further shows that pooled internal validation can overstate apparent clinical separation. This result should be presented as a strength rather than a weakness: the benchmark exposes the real generalization difficulty of GBM prognosis.

\textbf{The proposed model is promising as a mechanism, not yet as the top predictor.} The prior-guided kernel transformer recovered competitive 3-year performance but did not surpass the strongest classical baselines. The likely reasons are structural. First, the external branch currently omits methylation because TCGA probe harmonization is unfinished, so the model cannot yet exploit the modality that contributed most in the ablation study. Second, key clinical factors such as IDH mutation and MGMT promoter methylation are not yet integrated. Third, the scRNA prior bank is small. Recent glioma work combining EcoTyper-derived cell states and Self-Normalizing Networks reported much higher survival discrimination on larger mixed glioma cohorts \cite{fan2025tme}; our stricter GBM-only setting indicates that reproducing that level of performance requires both broader covariate coverage and more extensive cross-cohort harmonization.

\textbf{Practical next steps are now clear.} In the short term, the most valuable improvements are to harmonize TCGA methylation features for the external branch, add IDH/MGMT covariates, and report time-dependent survival metrics alongside fixed-horizon ROC. In the mid term, the scRNA branch should move from recurrent marker transfer to richer state-level scores or EcoTyper-style ecosystem recovery. In the longer term, the prior-guided transformer should become a truly multimodal cross-attention model in which RNA, methylation, and transferred cell-state features are encoded in separate streams rather than concatenated after the fact.
